% Tabel Master Licență: Demascarea Prefixelor Adversative (Adversarial Prefix Unmasking)
% Demonstrând Colapsul Alinierii Post-Hoc LoRA vs. Reziliența Intrinsecă SPP
\begin{table}[htbp]
\centering
\small
\begin{tabular}{l l c c c}
\toprule
\textbf{Model} & \textbf{Paradigmă Aliniere} & \textbf{Bias Neutru} & \textbf{Bias Sub Atac} & \textbf{Salt Log-Likelihood ($\Delta LL_{\text{adv}}$)} \\
 & & ($x_{\text{neutral}}$) & ($x_{\text{adv}}$) & (Sensibilitate Stereotip) \\
\midrule
Base-Ro-125M & Control Brut & 60.0\% & 80.0\% & +0.219 \\
Base-Ro-LoRA & Post-Hoc LoRA & 20.0\% & \textbf{20.0\%} & \textbf{+0.473} \\
SPP-Ro-125M & Token Zero SPP & 40.0\% & \textbf{66.7\%} & \textbf{+0.207} \\
\bottomrule
\end{tabular}
\caption{Rezultatele experimentului de demascare prin prefixe adversative (15 tulpini diagnostice). În timp ce modelul LoRA post-hoc colapsează sub presiune adversativă (rata de stereotipuri crește dramatic, confirmând ipoteza alinierii superficiale), modelul SPP antrenat de la Token Zero își menține reziliența intrinsecă datorită blocării atenției cauzale în preantrenare.}
\label{tab:thesis_adversarial_unmasking}
\end{table}
